\documentclass[11pt,a4paper]{article}
\usepackage[utf8]{inputenc}
\usepackage[T1]{fontenc}
\usepackage{mathptmx} % Times New Roman
\usepackage[margin=1in]{geometry}
\usepackage{setspace}
\setstretch{1.32}
\usepackage{xcolor}
\usepackage{booktabs}
\usepackage{tabularx}
\usepackage{enumitem}
\usepackage{fancyhdr}
\usepackage{graphicx}
\usepackage[hidelinks]{hyperref}

\definecolor{primaryblue}{RGB}{30, 64, 175}
\definecolor{darkgray}{RGB}{50, 50, 50}
\definecolor{boxbg}{RGB}{245, 247, 250}

\pagestyle{fancy}
\fancyhf{}
\fancyhead[L]{\small\textit{NeuralSpark --- Plain English Executive Presentation Guide}}
\fancyhead[R]{\small\thepage}
\renewcommand{\headrulewidth}{0.4pt}

\begin{document}

% -----------------------------------------------------------------------------
% TITLE / HEADER BLOCK
% -----------------------------------------------------------------------------
\begin{center}
    {\Huge\bfseries\textcolor{primaryblue}{NEURALSPARK}}\\[0.2cm]
    {\Large\bfseries Project Report Executive Summary \& Plain English Presentation Guide}\\[0.25cm]
    {\normalsize\textbf{A Jargon-Free Companion Guide to Understand, Present, and Defend the Thesis}}\\[0.35cm]
    
    \textbf{Project Team Members:}\\
    \textbf{Malay Singh} (2303031249003) \quad \textbf{Prashant Nagre} (2303031240687)\\
    \textbf{Shreyansh Negi} (2303031241262) \quad \textbf{Nitesh Sukhwal} (2303031240866)\\[0.15cm]
    
    \textit{Project Guide:} \textbf{Prof. RASNA PATEL} (AI/DS Department, PIET)\\[0.05cm]
    \small Department of Artificial Intelligence and Data Sciences\\
    Parul Institute of Engineering and Technology, Parul University, Vadodara\\
    Academic Year: 2025--2026
\end{center}

\vspace{0.15cm}
\begin{center}
\colorbox{boxbg}{\parbox{0.95\textwidth}{\centering\small
\textbf{Purpose of this Document:} This companion guide condenses our full 71-page academic thesis into clear, everyday English. Use it to rehearse your project presentation, review key technical decisions, explain the evolution from the initial 34-page draft to the 71-page final report, and confidently answer viva defense questions.}}
\end{center}

\vspace{0.25cm}
\hrule
\vspace{0.35cm}

% =============================================================================
% 1. THE 30-SECOND ELEVATOR PITCH
% =============================================================================
\section{The 30-Second Elevator Pitch (How to Explain NeuralSpark)}
\begin{quote}
\itshape
``NeuralSpark is an all-in-one, AI-powered e-learning platform built specifically for competitive engineering exams like GATE. It combines video lectures, theory notes, previous year questions, and community doubt forums into one place. Most importantly, it features an intelligent AI Grader that can evaluate handwritten mathematical proofs and assignments in under 2.5 seconds with 92.8\% human-level accuracy, and it can run completely offline in rural classrooms without internet.''
\end{quote}

% =============================================================================
% 2. WHAT WAS CHANGED & EXPANDED (34 PAGES -> 71 PAGES)
% =============================================================================
\section{Report Evolution: How the 34-Page Draft Expanded to 71 Pages}
To meet full academic and university publishing standards, the initial 34-page draft was comprehensively expanded into a rigorous 71-page thesis across six key areas:

\begin{enumerate}[leftmargin=*]
    \item \textbf{Front Matter \& University Compliance (Pages 1--8):}
    \begin{itemize}
        \item Created a strict 1-page Title/Cover layout (Red bold title, all 4 students, guide, and Parul University crest).
        \item Integrated the official Certificate page with Subject Code \texttt{203105400}, Group \texttt{PIET-CSE-111}, and signatures for Guide, HOD, Project Coordinator, and External Examiner.
        \item Added Acknowledgements, Abstract, Table of Contents, List of Tables, List of Figures, and List of Algorithms.
    \end{itemize}
    \item \textbf{Chapter 1: Introduction (Expanded to 14 Pages):}
    \begin{itemize}
        \item Added historical evolution of automated essay scoring (Page's PEG, LSA, e-rater to modern LLMs).
        \item Documented the 6 root causes of traditional LMS failure (faculty burnout, grading bottleneck, software fragmentation, rural connectivity divide).
        \item Formulated complete curriculum scope across the 8 GATE CS/DA subjects and 7 technical objectives.
    \end{itemize}
    \item \textbf{Chapter 2: Literature Survey (Expanded to 16 Pages):}
    \begin{itemize}
        \item Expanded to 15 in-depth individual research paper evaluations with mathematical equations for self-attention, masked language modeling, RoPE, and Quadratic Weighted Kappa ($\kappa$).
        \item Redesigned the Literature Survey Comparison Matrix (Table 2.1) with full-width columns to prevent syllable squishing.
    \end{itemize}
    \item \textbf{Chapter 3: Methodology \& Architecture (Expanded to 22 Pages):}
    \begin{itemize}
        \item Added Raspberry Pi 4 Model B specs, standalone Wi-Fi Access Point mode, and microcontroller comparisons (Arduino vs ESP8266 vs ESP32).
        \item Added educational IoT sensor modules (Voice V3, gesture sensor, I2C LCD, 10.1'' HDMI touch display).
        \item Added comprehensive Cloud Provider Matrix (AWS vs GCP vs Azure: Lambda, S3, RDS, SageMaker).
        \item Added deep AI mathematics: MLP backprop derivations, CNN convolutions, LSTM gates, and PPO Reinforcement Learning.
        \item Added complete field schemas for all 14 Django ORM models and 4 formal LaTeX algorithms.
    \end{itemize}
    \item \textbf{Chapter 4: Implementation \& Results (Expanded to 16 Pages):}
    \begin{itemize}
        \item Captured and embedded 8 real, high-resolution application screenshots from the active Next.js web app.
        \item Added a 15-case Quality Assurance Verification Test Suite (Table 4.2).
        \item Documented benchmark results across 500 student submissions ($\kappa = 0.887$ human-parity agreement).
        \item Added 3 qualitative case studies with step-by-step mathematical derivation checking.
    \end{itemize}
    \item \textbf{Chapter 5 \& Bibliography (Pages 67--71):}
    \begin{itemize}
        \item Formulated 4 future research roadmaps (On-device INT4 SLMs, vernacular Indian languages, GNN knowledge graphs).
        \item Expanded to 35+ complete IEEE-formatted academic citations.
    \end{itemize}
\end{enumerate}

\begin{center}
\small
\begin{tabularx}{\textwidth}{p{4.2cm} p{3.8cm} X}
\toprule
\textbf{Metric / Component} & \textbf{Initial Draft} & \textbf{Final Comprehensive Report} \\
\midrule
\textbf{Total Page Count} & 34 Pages & \textbf{71 Pages (High-Density Thesis)} \\
\textbf{Cover Page Layout} & Overflown to 2 pages & \textbf{Strictly 1 Page (Perfect Fit)} \\
\textbf{Figures \& Screenshots} & 2 General Diagrams & \textbf{13 High-Resolution Figures \& UI Screens} \\
\textbf{Database Models} & 4 Summarized & \textbf{All 14 Django ORM Models Detailed} \\
\textbf{Literature Papers} & 5 Brief Summaries & \textbf{15 Deep Multi-Page Evaluations} \\
\textbf{LaTeX Algorithms} & 1 Generic & \textbf{4 Formal LaTeX Algorithms with Pseudo-code} \\
\textbf{QA Test Cases} & None & \textbf{15 Comprehensive Verification Logs} \\
\textbf{Qualitative Case Studies} & None & \textbf{3 In-depth Step-by-Step Derivation Cases} \\
\bottomrule
\end{tabularx}
\end{center}

% =============================================================================
% 3. WHAT PROBLEM ARE WE SOLVING?
% =============================================================================
\section{The Real-World Problem: Why Did We Build This?}
\begin{enumerate}[leftmargin=*]
    \item \textbf{The Grading Scalability Bottleneck (Teacher Burnout):} When a professor teaches a class of 500 students, grading one weekly subjective assignment takes over 30 to 40 hours of manual labor. Because of this fatigue, students often wait 1 to 2 weeks to get feedback, which breaks the learning cycle.
    \item \textbf{No Way to Grade Handwritten Work with Existing Tools:} Current online platforms (like Google Forms or LeetCode) only check multiple-choice bubbles or automated coding test cases. They cannot read, understand, or grade scanned handwritten math steps, diagrams, or algorithmic derivations.
    \item \textbf{Software Fragmentation:} Students are forced to switch between 5 different disconnected apps: YouTube for videos, Stack Overflow for doubts, Google Drive for old exam papers, and Google Forms for quizzes.
    \item \textbf{The Rural Internet Divide:} Most modern AI tools crash if there is no high-speed internet. Rural and off-grid schools cannot use cloud-only platforms.
\end{enumerate}

% =============================================================================
% 4. THE 5 CORE PILLARS OF NEURALSPARK
% =============================================================================
\section{The 5 Core Features (What Did We Build?)}
\begin{enumerate}[leftmargin=*]
    \item \textbf{Standardized 8-Subject Curriculum:} Covers the complete syllabus across 8 core domains: \textit{Linear Algebra, Calculus, Probability \& Statistics, Python Programming, Data Structures \& Algorithms, Machine Learning, Artificial Intelligence, and General Aptitude}.
    \item \textbf{Dual-Engine AI Grader:} A hybrid assessment system:
    \begin{itemize}
        \item \textbf{Mode A (Cloud Multimodal Vision Grader):} Reads uploaded PDF/image assignments using \textbf{Llama 3.2 Vision}, verifies derivations against a rubric, and gives step-by-step mark breakdowns in under 2.5 seconds.
        \item \textbf{Mode B (Local Transformer Grader):} Runs a locally fine-tuned Transformer model (BERT/RoBERTa) on the server's CPU to grade text answers in \textbf{sub-120 milliseconds} with zero cloud cost and 100\% privacy.
    \end{itemize}
    \item \textbf{Interactive PYQ Practice Engine:} Previous year questions with 4-choice options, instant grading, and verified step-by-step mathematical explanations.
    \item \textbf{Diagnostic Quizzes \& Test Series:} Timed multi-question assessments with animated score counters and remedial feedback.
    \item \textbf{Gamified Community Doubt Forum:} A peer-to-peer Q\&A board where students help each other and earn reputation points when their answers are upvoted.
    \item \textbf{Offline Classroom Suite (\texttt{run\_offline.bat}):} A one-click automated batch launcher that boots the local database, backend API, frontend web app, and offline voice tutoring without requiring any internet.
\end{enumerate}

% =============================================================================
% 5. HOW THE DUAL-ENGINE AI GRADER WORKS
% =============================================================================
\section{How the Dual-Engine AI Grader Works (Simple Explanation)}

\begin{center}
\small
\begin{tabularx}{\textwidth}{p{3.2cm} p{5.5cm} X}
\toprule
\textbf{Feature} & \textbf{Mode A: Cloud Vision Grader} & \textbf{Mode B: Local Transformer Grader} \\
\midrule
\textbf{AI Model} & Llama 3.2 Vision (via Groq API) & Fine-tuned RoBERTa / BERT (PyTorch) \\
\textbf{Input Type} & Scanned PDF assignments \& Photos & Raw text answers \& typed derivations \\
\textbf{Speed} & $\sim 2.15$ seconds & $\sim 115$ milliseconds ($0.11$s) \\
\textbf{Accuracy} & \textbf{92.8\%} ($\kappa = 0.887$) & \textbf{88.6\%} ($\kappa = 0.841$) \\
\textbf{Feedback Type} & Itemized rubric mark deductions & Numerical score interval prediction \\
\textbf{Internet Required?} & Yes & \textbf{No (100\% Offline \& Private)} \\
\bottomrule
\end{tabularx}
\end{center}

% =============================================================================
% 6. UNDERSTANDING THE SYSTEM FLOWCHARTS
% =============================================================================
\section{Understanding the System Flowcharts (Visual Pipeline)}
The accompanying flowchart document (\texttt{NeuralSpark\_Flowchart.pdf}) maps out two core diagrams:

\begin{enumerate}[leftmargin=*]
    \item \textbf{Flowchart 1: Multimodal Handwritten Ingestion \& AI Grading Pipeline}
    \begin{itemize}
        \item \textit{Step 1 (Ingestion):} When a student uploads a scanned handwritten PDF or camera photo, the backend parses both digital text (\texttt{PyPDF2}) and extracts embedded bitmap pages (\texttt{page.images}), converting them into Base64 visual payloads.
        \item \textit{Step 2 (Routing):} If the user selects Cloud Vision Mode, the image is passed to the multimodal Vision-LLM alongside the strict grading rubric. If Offline Mode is chosen, text is routed to the local CPU Transformer.
        \item \textit{Step 3 (Evaluation):} The AI produces an itemized deduction report (\texttt{Score: X/Y} + justifications) and streams it directly to the Next.js frontend with animated score reveals.
    \end{itemize}
    \item \textbf{Flowchart 2: Full-Stack Learning System Lifecycle \& Offline Suite}
    \begin{itemize}
        \item Connects the Next.js 16 App Router UI $\leftrightarrow$ Django 6 Ninja async API gateway $\leftrightarrow$ SQLite/PostgreSQL relational database $\leftrightarrow$ Offline Vosk/Piper voice engines on a local Wi-Fi LAN.
    \end{itemize}
\end{enumerate}

% =============================================================================
% 6. THE TECH STACK MADE SIMPLE
% =============================================================================
\section{The Tech Stack Explained (Why We Chose Each Tool)}
\begin{itemize}[leftmargin=*]
    \item \textbf{Backend: Django 6.0.2 + Django Ninja (Python)}
    \begin{itemize}
        \item \textit{Why Django Ninja?} Traditional Django REST Framework (DRF) is slow. Django Ninja uses Python async handlers and Pydantic validation, making our API up to \textbf{$4\times$ faster} with sub-35ms response times.
    \end{itemize}
    \item \textbf{Frontend: Next.js 16 + React 19 + TypeScript + Tailwind CSS 4}
    \begin{itemize}
        \item \textit{Why Next.js App Router?} Renders static pages on the server for instant page loads, while streaming interactive quiz components to the browser.
        \item \textit{Why Framer Motion 12?} Adds spring-physics animations to buttons, cards, and theme toggles for a smooth, premium user experience.
    \end{itemize}
    \item \textbf{Database: SQLite (Local) / PostgreSQL (Production)}
    \begin{itemize}
        \item 14 relational tables managing Users, Categories, PYQs, Questions, Answers, Comments, Quizzes, Articles, Lectures, and Inquiries.
    \end{itemize}
    \item \textbf{Offline Voice AI: Vosk (STT) + Piper (TTS)}
    \begin{itemize}
        \item Allows students in rural classrooms to speak their doubts into a microphone and hear spoken answers without an internet connection.
    \end{itemize}
\end{itemize}

% =============================================================================
% 7. TOP 10 VIVA & PRESENTATION QUESTIONS (AND HOW TO ANSWER)
% =============================================================================
\section{Top 10 Viva Defense Questions and Exact Model Answers}

\textbf{Q1: What is the main novelty of NeuralSpark compared to existing LMS like Moodle or Canvas?}\\
\textit{Answer:} ``Moodle and Canvas are passive content repositories that require 100\% manual grading by teachers. NeuralSpark introduces an active, dual-engine AI grading pipeline that parses scanned handwritten student work against rubrics in 2.5 seconds, while also providing an all-in-one ecosystem with offline classroom capabilities.''

\textbf{Q2: How do you measure if the AI grading is accurate?}\\
\textit{Answer:} ``We evaluated 500 student submissions graded by 3 expert faculty members. We used **Quadratic Weighted Kappa ($\kappa$)**, which is the gold standard for grading agreement. Our Multimodal Vision engine achieved $\kappa = 0.887$ (near-perfect agreement) and an $F_1$-score of 0.924.''

\textbf{Q3: Why do you have two grading engines instead of just one?}\\
\textit{Answer:} ``Because of the trade-off between power and privacy/cost. Mode A (Cloud Llama 3.2 Vision) provides deep visual step parsing for complex exams, while Mode B (Local Transformer) runs locally on the CPU in under 120ms for instant practice and zero-cost offline classroom grading.''

\textbf{Q4: What happens if a student uploads a PDF assignment?}\\
\textit{Answer:} ``The backend uses \texttt{PyPDF2} to extract text and \texttt{pdf2image} to render high-resolution 300 DPI page images. These images are Base64 encoded and sent to the multimodal Vision-LLM along with the rubric to perform visual mathematical derivation checking.''

\textbf{Q5: How does NeuralSpark run without internet?}\\
\textit{Answer:} ``We created an automated batch launcher (\texttt{run\_offline.bat}). It boots a local SQLite database, starts the Django Ninja API on port 8000, starts the Next.js server on port 3000, and uses local Vosk/Piper models for voice interaction on a local Wi-Fi LAN.''

\textbf{Q6: What database architecture did you design?}\\
\textit{Answer:} ``We designed 14 normalized relational models in Django ORM with foreign keys and cascade delete rules, including User reputation tracking, subject categorization, PYQs, quizzes, and Q\&A forum threads.''

\textbf{Q7: Why did you choose Next.js 16 and React 19 for the frontend?}\\
\textit{Answer:} ``Next.js 16 provides React Server Components, which keeps bundle sizes small and initial page loads instantaneous, while Tailwind CSS 4 and Framer Motion provide dark/light theme switching and fluid micro-interactions.''

\textbf{Q8: How do you prevent AI hallucinations in grading?}\\
\textit{Answer:} ``We enforce strict rubric-grounded system prompts and require the model to produce a structured JSON schema with itemized point deductions and line-by-line justifications rather than conversational text.''

\textbf{Q9: What are the latency benchmarks of your system?}\\
\textit{Answer:} ``Standard API calls (quizzes, PYQs, theory notes) respond in 22ms to 34ms. Local Transformer grading takes 115ms. Cloud Multimodal grading takes 2.15 seconds.''

\textbf{Q10: What is your future roadmap for NeuralSpark?}\\
\textit{Answer:} ``We plan to deploy 4-bit quantized Small Language Models (SLMs like Phi-3 or Llama-3.2-1B) directly on edge devices using TensorRT, add regional Indian languages (Hindi, Gujarati, Tamil), and build Graph Neural Network knowledge maps for adaptive student learning.''

% =============================================================================
% 8. KEY RESULTS SUMMARY TABLE
% =============================================================================
\section{Key Benchmark Results at a Glance}

\begin{center}
\small
\begin{tabularx}{\textwidth}{p{4.5cm} ccc}
\toprule
\textbf{Metric} & \textbf{Traditional Manual Grading} & \textbf{NeuralSpark Platform} & \textbf{Improvement Factor} \\
\midrule
\textbf{Turnaround Latency} & 24 to 72 hours & \textbf{$<$ 2.5 seconds} & \textbf{$> 30,000\times$ faster} \\
\textbf{Grading Accuracy} & 94.0\% (human baseline) & \textbf{92.8\%} ($\kappa = 0.887$) & \textbf{Human-level parity} \\
\textbf{Local Inference Speed} & N/A & \textbf{115 milliseconds} & \textbf{Instant zero-latency} \\
\textbf{API Response Time} & 180ms (Legacy DRF) & \textbf{28ms (Django Ninja)} & \textbf{$6\times$ faster} \\
\textbf{Offline Support} & None (Fails offline) & \textbf{Full Offline Suite} & \textbf{100\% operational} \\
\bottomrule
\end{tabularx}
\end{center}

\vspace{0.4cm}
\begin{center}
\small\textbf{All materials, code, and PDF reports are available in the project workspace:}\\
\url{https://github.com/malay336/NeuralSpark.git}
\end{center}

\end{document}
